\section*{РЕФЕРАТ}


Пояснювальна записка до дипломної роботи <<Метод масштабування цифрових зображень на основі другої барицентричної інтерполяційної формули за вузлами Чебишева>>: 69~сторінок, 14~рисунків, 6~використаних джерел.

\medskip

МАСШТАБУВАННЯ ЗОБРАЖЕНЬ, БАРИЦЕНТРИЧНА ІНТЕРПОЛЯЦІЯ, ВУЗЛИ ЧЕБИШЕВА ДРУГОГО РОДУ, ФЕНОМЕН РУНГЕ, ФЕНОМЕН ГІББСА, ЦИФРОВА ОБРОБКА ЗОБРАЖЕНЬ, МАТРИЧНА ФАКТОРИЗАЦІЯ, C++, QT, OPENCV, PSNR, SSIM

\medskip

\textbf{Об'єкт дослідження} --- обробка цифрових зображень. \textbf{Предмет дослідження} --- метод масштабування цифрових зображень за допомогою другої барицентричної формули по вузлах Чебишева другого роду.

\textbf{Мета роботи} --- розробити та програмно реалізувати новий метод масштабування цифрових зображень на основі другої барицентричної інтерполяційної формули на вузлах Чебишева другого роду.

\textbf{Методи дослідження:} математичне моделювання на основі другої барицентричної інтерполяційної формули, апарат лінійної алгебри та засоби розпаралелювання обчислень, кількісне оцінювання якості масштабування за метриками PSNR та SSIM. Технічні та програмні засоби: мова програмування C++, фреймворк Qt для побудови графічного інтерфейсу, бібліотека комп'ютерного зору OpenCV, середовище MATLAB для генерації еталонних зображень.

У роботі математичну модель цифрового зображення адаптовано до нерівномірних сіток дискретизації за вузлами Чебишева--Лобатто, що усуває феномен Рунге та забезпечує властивість вкладеності сіток. Другу барицентричну формулу поширено на двовимірний випадок; шляхом факторизації інтерполяційного ядра подвійну суму зведено до добутку трьох матриць ($G = B \cdot F \cdot A^{\top}$), що знизило асимптотичну складність глобальної інтерполяції до послідовної одновимірної. Для придушення феномена Гіббса на різких контурах запропоновано локальну модифікацію методу на ковзному вікні $K \times K$. Створено програмний комплекс із графічним інтерфейсом, підтримкою обробки виділених областей (ROI) та модулем автоматизованого тестування з порівнянням розробленого методу з індустріальними алгоритмами OpenCV (найближчого сусіда, білінійним, бікубічним).

\textbf{Наукова новизна} полягає в розвитку математичної моделі представлення цифрового зображення на сітці вузлів Чебишева--Лобатто, виведенні та обґрунтуванні двовимірної форми другої барицентричної формули, а також оптимізації обчислювального процесу шляхом факторизації двовимірного ядра та побудови локальної барицентричної інтерполяції на обмеженому просторовому околі.

\textbf{Значущість роботи} полягає у доведенні теоретичних положень до рівня конкретного прикладного рішення --- математично строгого, чисельно стійкого та апаратно-прискореного алгоритму масштабування, здатного працювати в режимі реального часу. Результати автоматизованого тестування за метриками PSNR та SSIM підтвердили конкурентоспроможність методу порівняно з класичними алгоритмами.

Результати рекомендується використовувати у спеціалізованих задачах, де вимоги до відсутності артефактів та структурної точності даних є критично вищими за вимоги до швидкості обчислень. \textbf{Галузь застосування} --- комп'ютерний зір, медична візуалізація, аналіз супутникових знімків, астрофізика, системи відеоспостереження. Основні положення роботи апробовано на Міжнародній науково-практичній конференції <<Політ--2026>>.

Подальший розвиток предмета дослідження передбачає перенесення матричних обчислень на графічні прискорювачі (CUDA, OpenCL) для обробки растрових масивів надвисокої роздільної здатності та відеопотоків у реальному часі, а також застосування сум і фільтрів де ла Валле-Пуссена для додаткового придушення залишкових осциляцій без втрати різкості зображення.